\section{Results}
\label{sec:results}

Present your findings objectively. Use tables and figures to support the text.
Describe what is shown without yet interpreting it (save interpretation for
\cref{sec:conclusion}).

\subsection{Quantitative Results}

\Cref{tab:results} shows the performance of each method.

\begin{table}[h]
  \centering
  \caption{Comparison of methods on the test set.}
  \label{tab:results}
  \begin{tabular}{lrrc}
    \toprule
    Method       & Metric A (\%) & Metric B (s) & Notes   \\
    \midrule
    Baseline     & 72.3          & 1.45         & —       \\
    Ours (v1)    & 85.1          & 0.92         & +12.8\% \\
    Ours (v2)    & \textbf{91.4} & \textbf{0.78} & +19.1\% \\
    \bottomrule
  \end{tabular}
\end{table}

\subsection{Visual Analysis}

\begin{figure}[h]
  \centering
  \begin{tikzpicture}
    \begin{axis}[
      width=0.72\textwidth, height=5cm,
      ybar, bar width=18pt,
      symbolic x coords={Baseline, Ours v1, Ours v2},
      xtick=data,
      ylabel={Metric A (\%)},
      ymin=60, ymax=100,
      nodes near coords,
      every node near coord/.style={font=\small},
      enlarge x limits=0.3,
    ]
      \addplot[fill=codeblue!70] coordinates {
        (Baseline,72.3) (Ours v1,85.1) (Ours v2,91.4)
      };
    \end{axis}
  \end{tikzpicture}
  \caption{Metric A scores across methods.}
  \label{fig:results-bar}
\end{figure}

\subsection{Code Snapshot}

\begin{lstlisting}[language=Python, caption={Core algorithm implementation.}]
def run_experiment(config):
    model = build_model(config)
    train(model, config.train_data)
    metrics = evaluate(model, config.test_data)
    return metrics
\end{lstlisting}
